\chapter{La fibre longue distance}\label{ch:nw:long-haul-fibre}

En août 2010, une nouvelle route de fibre entre Chicago et le New Jersey entra en service~: plus courte que les routes qui l'avaient précédée, avec un équipement à plus faible latence,
et annoncée à \qty{6.65}{\milli\second} à sens unique entre le bâtiment Cermak de Chicago et le centre de données de Nasdaq à Carteret,
soit environ \qty{13.3}{\milli\second} aller-retour. Son coût fut estimé publiquement à environ 300~millions de dollars. Avant elle, les cours du New Jersey
réagissaient à Chicago en 7.25 à \qty{7.95}{\milli\second}~; autour de son ouverture, les temps de réaction tombent de 1 à
\qty{1.5}{\milli\second}. En deux ans, les réseaux hertziens
estimés à 4.2 à \qty{5.2}{\milli\second} avaient fait mieux qu'elle, et les acheteurs du câble avaient appris ce que disent les planchers du chapitre~10~: le verre est
plus lent que l'air de 46~\%, et aucune quantité de travaux n'y change rien.

La fibre n'a pas perdu. Elle transporte tout ce qui n'est pas parmi les quelques messages les plus rapides~: les données de marché au gigabit, les ordres lorsque la météo coupe
les liaisons radio (chapitre~14), les routes transatlantiques que nulle tour ne peut franchir, et chaque connexion entre l'armoire d'une société et son opérateur.
Ce chapitre démonte une route longue distance~: le verre et ce qui arrive à la lumière qui le traverse, le tracé et sa rectitude, l'équipement, les
câbles sous-marins, les façons d'acheter de la fibre, et le nouveau verre, presque aussi rapide que l'air.

\section{Lumière dans le verre~: atténuation, amplificateurs et dispersion}

La fibre monomode standard à \qty{1550}{\nano\metre} perd environ \qty{0.16}{\decibel} par kilomètre (la fiche technique de Corning pour sa fibre à très faible atténuation
donne 0.16 en valeur typique et 0.17 au maximum), si bien qu'un signal perd la moitié de sa puissance tous les 19~kilomètres et un
millième tous les 190. Sur les \qty{1359}{\kilo\metre} de verre que la route de 2010 doit contenir (voir plus bas), la perte serait de
\qty{217}{\decibel}~: aucun récepteur ne capte cela. La fibre longue distance est donc découpée en tronçons, avec un amplificateur à l'extrémité de chacun.

\begin{definition}[Amplificateur optique]\label{def:nw:long-haul-fibre:amp}
Un \emph{amplificateur optique}\index{amplificateur optique} restaure la puissance de la lumière dans une fibre sans la convertir en signal électrique,
généralement en la faisant passer par quelques mètres de fibre dopée à une terre rare et pompée par un laser~; placé à l'extrémité de chaque tronçon,
il permet à un signal de traverser des milliers de kilomètres de verre.
\end{definition}

L'amplificateur ne coûte guère de temps (la fibre dopée fait quelques mètres, le pass-through une fraction de microseconde), raison pour laquelle le modèle de route ci-dessous
compte \qty{0.1}{\micro\second} par tronçon, une hypothèse, et pourquoi l'amplification n'est pas le problème de latence. C'est la dispersion.

\begin{definition}[Dispersion chromatique, compensation de dispersion]\label{def:nw:long-haul-fibre:disp}
La \emph{dispersion chromatique}\index{dispersion chromatique} est l'étalement d'une impulsion lumineuse dans une fibre, car ses différentes longueurs d'onde se propagent
à des vitesses légèrement différentes, mesuré en picosecondes d'étalement par nanomètre de largeur spectrale et par kilomètre. La \emph{compensation
de dispersion}\index{compensation de dispersion} la défait, soit optiquement, avec une longueur de fibre de dispersion opposée (fibre compensatrice
de dispersion) dans chaque tronçon, soit électroniquement, par traitement du signal dans un récepteur cohérent.
\end{definition}

La même fiche technique borne la dispersion à \qty{18}{\pico\second} par nanomètre et par kilomètre à \qty{1550}{\nano\metre}~: sur 1\,359~kilomètres, l'étalement atteint environ
\qty{24}{\nano\second} par nanomètre de spectre, bien plus qu'un bit à dix gigabits (\qty{0.1}{\nano\second}). La compenser
optiquement ajoute du verre au trajet~: une bobine dans chaque tronçon, que le modèle fixe au cinquième de la longueur du tronçon. La compenser électroniquement
ajoute en revanche du temps de traitement aux extrémités, et laisse le trajet intact.

\begin{definition}[Correction d'erreurs directe]\label{def:nw:long-haul-fibre:fec}
La \emph{correction d'erreurs directe}\index{correction d'erreurs directe} (FEC) ajoute des bits redondants à un signal transmis, afin que le récepteur puisse
corriger les erreurs sans demander de retransmission~; son décodeur doit rassembler un bloc (ou un groupe entrelacé de blocs) avant de pouvoir
la corriger, ce qui ajoute un délai à l'extrémité de réception.
\end{definition}

La FEC permet à une liaison de fonctionner avec un signal plus bruité, donc des tronçons plus longs ou des débits plus élevés, au prix d'un délai par terminal. Laughlin, Aguirre et
Grundfest signalent le même compromis dans les faisceaux hertziens, dont la latence est «~dominée par~» la mise en mémoire tampon et l'entrelacement de la FEC. Pour la fibre à faible latence,
les choix sont donc ceux que le modèle expose~: tronçons courts ou longs, \omterm{def:nw:long-haul-fibre:disp}{compensation de dispersion} optique ou électronique, FEC lourde
ou légère.

\begin{omfigure}
\begin{tikzpicture}[font=\footnotesize,b/.style={draw,rounded corners=2pt,minimum height=0.6cm,align=center,inner sep=3pt},>=Stealth]
\node[b,fill=omThm!15] (t1) at (0,0) {terminal~:\\ transpondeur\\ et FEC};
\node[b,fill=omDef!15] (a1) at (4.6,0) {amplificateur};
\node[b,fill=omDef!15] (a2) at (8.4,0) {amplificateur};
\node[b,fill=omThm!15] (t2) at (12.4,0) {terminal~:\\ décodage FEC\\ et transpondeur};
\draw[thick] (t1) -- node[above=2pt]{tronçon 1~: 80~km} (a1);
\draw[thick] (a1) -- node[above]{tronçon 2} (a2);
\draw[thick,dashed] (a2) -- node[above]{\dots\ tronçon $N$} (t2);
\draw[omMeth,thick] (6.5,-0.95) circle (0.28) node[below=0.3cm,font=\scriptsize,omMeth]{bobine DCF optionnelle};
\draw[omMeth,->] (6.5,-0.67) -- (6.5,-0.05);
\node[font=\scriptsize,align=center] at (6.2,1.1) {latence = verre le long de la géodésique + trajet supplémentaire + DCF + amplificateurs + terminaux};
\end{tikzpicture}

\omcaption{Une route de fibre longue distance, schématiquement~: terminaux à chaque extrémité (transpondeurs et \omterm{def:nw:long-haul-fibre:fec}{correction d'erreurs directe}), tronçons de verre avec un
amplificateur à l'extrémité de chacun, et, sur les routes qui compensent la dispersion optiquement, une bobine de fibre compensatrice de dispersion (DCF) dans chaque tronçon.
\texttt{firm.fibreroute} attribue à chaque élément son temps.}\label{fig:nw:long-haul-fibre:span}
\end{omfigure}

\section{Rectitude du tracé~: l'histoire Chicago--New York}

Le plancher du chapitre~10 pour 350 Cermak à Carteret vaut \qty{5507}{\micro\second} en fibre droite~; la route de 2010 était annoncée à
\qty{6650}{\micro\second}. L'écart, \qty{1143}{\micro\second}, tient au tracé et à l'équipement, et inverser le chiffre publié les sépare
une fois un temps d'équipement posé. Avec deux terminaux de \qty{10}{\micro\second} et un amplificateur tous les 80~kilomètres à
\qty{0.1}{\micro\second} (hypothèses du modèle, \qty{21.7}{\micro\second} au total), la route doit contenir 1\,359~kilomètres de verre~: 230
kilomètres de plus que la géodésique, soit un \omterm{def:nw:the-north-american-map:factor}{facteur de tracé} de 1.20. Le chiffre ultérieur de \qty{6.55}{\milli\second} implique 1\,339~kilomètres, 20 de moins,
ou un équipement plus rapide.

\begin{proposition}[Inversion d'une latence publiée]\label{prop:nw:long-haul-fibre:invert}
Si une route de fibre d'\omterm{def:nw:the-north-american-map:floor}{indice de groupe} $n_g$ est annoncée à $T$ à sens unique et que son équipement prend $E$, son verre vaut
\[
  L = \frac{(T - E)\,c_0}{n_g},
\]
et son \omterm{def:nw:the-north-american-map:factor}{facteur de tracé} vaut $L/d$ pour une géodésique $d$. Une erreur $\delta E$ sur le temps d'équipement supposé déplace $L$ de $\delta E\,c_0/n_g$, soit environ
\qty{205}{\metre} par microseconde.
\end{proposition}

\begin{proof}
La lumière parcourt $L$ à $c_0/n_g$ dans le temps qui reste une fois le $E$ de l'équipement retiré. En dérivant par rapport à $E$, on obtient la sensibilité.
\end{proof}

\omcode{code/firm/fibreroute/firm_fibreroute.py}{59}{75}{Latence d'une route par composante, et inversion d'une latence publiée en
verre que la route doit contenir.}\label{lst:nw:long-haul-fibre:components}

\begin{omfigure}
\begin{tikzpicture}
\begin{axis}[width=11cm,height=5.6cm,xmin=0,xmax=200,ymin=1.16,ymax=1.215,xlabel={temps d'équipement supposé (\si{\micro\second}, à sens unique)},
  ylabel={facteur de tracé implicite},tick label style={font=\footnotesize},label style={font=\small},grid=major,grid style={gray!25},
  table/col sep=comma,yticklabel style={/pgf/number format/fixed,/pgf/number format/precision=2}]
\addplot[omDef,thick,no markers] table[x=equipment_us,y=factor]{figdata/networks/13-long-haul-fibre/sensitivity.csv};
\end{axis}
\end{tikzpicture}

\omcaption{Quelle part de l'excédent de la route de 2010 relève du tracé~: le \omterm{def:nw:the-north-american-map:factor}{facteur de tracé} implicite par sa latence annoncée de \qty{6.65}{\milli\second}, en fonction du
temps d'équipement supposé. Même avec \qty{200}{\micro\second} d'équipement, un trajet reste supérieur de 17~\% à la géodésique. Données~: \texttt{fig\_fibre.py},
\texttt{nw\_fibre.sensitivity()}.}\label{fig:nw:long-haul-fibre:sensitivity}
\end{omfigure}

La \Cref{fig:nw:long-haul-fibre:sensitivity} montre pourquoi la conclusion est robuste~: quel que soit l'équipement, l'essentiel de la milliseconde relève du tracé. Un câble
doit suivre les droits de passage qu'il peut obtenir~; la route de 2010 a acheté la rectitude pour 300~millions de dollars et restait pourtant plus longue de 17 à 20~\% que la
géodésique. Les réseaux hertziens qui la devançaient avaient des \omterm{def:nw:the-north-american-map:factor}{facteurs de tracé} proches de 1.01 dans l'air (chapitre~10) et un milieu de propagation 46~\% plus rapide.

\begin{table}[ht]
\centering\small
\begin{tabular}{lrrrrr}
\toprule
 & (km) & (ms, sens unique) & (ms) & (km) & facteur \\
 & (km) & (ms, sens unique) & (ms) & (km) & facteur \\
\midrule
Chicago--New Jersey, 2010 & 1\,129.2 & 6.650 & 5.507 & 1\,359 & 1.204 \\
Chicago--New Jersey, plus tard & 1\,129.2 & 6.550 & 5.507 & 1\,339 & 1.185 \\
New York--London, sous-marin, 2015 & 5\,551.0 & 29.475 & 27.070 & 6\,038 & 1.088 \\
\bottomrule
\end{tabular}
\caption{Routes à fibre publiées, ramenées en longueur de verre et en \omterm{def:nw:the-north-american-map:factor}{facteur de tracé} (\cref{prop:nw:long-haul-fibre:invert}), avec les hypothèses d'équipement du modèle
hypothèses (\qty{10}{\micro\second} par terminal, \qty{0.1}{\micro\second} par tronçon de 80~km)~: les chiffres cités par Laughlin, Aguirre et Grundfest
pour la route de Chicago, et la moitié de l'aller-retour publié par Hibernia Express pour l'Atlantique. Données~: \texttt{nw\_fibre.published()}.}\label{tab:nw:long-haul-fibre:published}
\end{table}

\section{Câbles sous-marins}

\begin{definition}[Station d'atterrissement de câble]\label{def:nw:long-haul-fibre:cls}
Une \emph{station d'atterrissement de câble}\index{station d'atterrissement de câble} (cable landing station) est le bâtiment où un câble sous-marin accoste et se raccorde aux
réseaux terrestres~: elle alimente les répéteurs sous-marins du câble, termine ses paires de fibres et confie leur trafic à des liaisons à terre.
\end{definition}

Une route sous-marine est une route terrestre à chaque extrémité, une station d'atterrissement sur chaque côte, et l'océan entre les deux. Hibernia Express, le premier nouveau
câble transatlantique depuis plus de douze ans à son ouverture en 2015, a publié un aller-retour inférieur à \qty{58.95}{\milli\second} entre le NY4 d'Equinix
à Secaucus et le LD4 à Slough, soit plus de cinq millisecondes de moins que les câbles existants, sur une paire de fibres de Corning à cœur de silice pure
La géodésique entre les deux bâtiments mesure \qty{5551}{\kilo\metre}~; la moitié de l'aller-retour, \qty{29.475}{\milli\second}, implique environ
6\,040~kilomètres de verre, soit un \omterm{def:nw:the-north-american-map:factor}{facteur de tracé} de 1.09. Sur un océan, la géodésique est presque entièrement accessible, hormis les atterrissages et les raccordements
vers les centres de données~; sur terre, les droits de passage coûtent à la route de Chicago deux fois plus, proportionnellement.

\begin{remark}[Allers-retours et moitiés]\label{rem:nw:long-haul-fibre:half}
Diviser par deux un aller-retour publié suppose que les deux sens empruntent le même tracé et le même équipement. Sur une paire de fibres d'un câble, c'est généralement le cas~; sur un réseau dont
les deux sens empruntent des fibres ou des routes distinctes, ce n'est pas forcément le cas, et l'asymétrie du chapitre~4 revient~: un aller-retour ne permet pas de distinguer les deux
sens.
\end{remark}

\section{Fibre noire contre longueurs d'onde louées}

\begin{definition}[Fibre noire, service allumé, service de longueur d'onde]\label{def:nw:long-haul-fibre:dark}
Une \emph{fibre noire}\index{fibre noire} (dark fibre) est une fibre vendue ou louée sans équipement~: l'acheteur l'allume avec ses propres transpondeurs et amplificateurs et
contrôle tout ce qui affecte la latence. Un \emph{service allumé}\index{service allumé} (lit service) est une capacité vendue sur l'équipement du vendeur~: un circuit d'un
débit donné entre deux points. Un \emph{service de longueur d'onde}\index{service de longueur d'onde} (wavelength service) est un service allumé qui réserve une longueur d'onde du
système de fibres du vendeur à l'acheteur.
\end{definition}

\begin{definition}[Droit irrévocable d'usage]\label{def:nw:long-haul-fibre:iru}
Un \emph{droit irrévocable d'usage}\index{droit irrévocable d'usage} (indefeasible right of use, IRU) est un contrat de longue durée conférant à son titulaire le droit exclusif d'utiliser
des fibres ou une capacité désignées d'un réseau pour une durée fixe, généralement payée d'avance, sans transfert de propriété du câble, et
avec une part de la maintenance à la charge du titulaire.
\end{definition}

Le rapport annuel d'un opérateur internet montre la forme en pratique~: l'essentiel de sa \omterm{def:nw:long-haul-fibre:dark}{fibre noire} est détenu «~sous forme de baux de longue durée
au titre de \omterm{def:nw:long-haul-fibre:iru}{droits irrévocables d'usage}~», il s'appuie sur le propriétaire de la fibre pour l'entretenir et il amortit ses IRU sur 15 à 20 ans en général. Pour une
société de trading, le choix entre \omterm{def:nw:long-haul-fibre:dark}{fibre noire} et \omterm{def:nw:long-haul-fibre:dark}{service allumé} revient à arbitrer entre contrôle et coût. La \omterm{def:nw:long-haul-fibre:dark}{fibre noire} permet à la société de choisir son propre
équipement, avec une \omterm{def:nw:long-haul-fibre:fec}{correction d'erreurs directe} (FEC) allégée, une compensation électronique de dispersion et aucun commutateur sur le trajet, et coûte un IRU, l'équipement et la maintenance annuelle~; un
\omterm{def:nw:long-haul-fibre:dark}{service allumé} coûte des frais mensuels et donne la latence du vendeur.

\begin{omfigure}
\begin{tikzpicture}
\begin{axis}[width=11cm,height=5.8cm,xmin=0,xmax=20,ymin=0,ymax=23,xlabel={années},ylabel={coût cumulé (millions de dollars)},
  tick label style={font=\footnotesize},label style={font=\small},grid=major,grid style={gray!25},
  legend style={font=\footnotesize,cells={anchor=west},at={(0.5,-0.22)},anchor=north,draw=none},table/col sep=comma]
\addplot[omDef,thick,no markers] table[x=year,y=dark_m]{figdata/networks/13-long-haul-fibre/costs.csv};
\addplot[omThm,thick,dashed,no markers] table[x=year,y=lit_m]{figdata/networks/13-long-haul-fibre/costs.csv};
\legend{{fibre noire (IRU et équipement d'avance, maintenance annuelle)},longueur d'onde allumée (frais mensuels)}
\end{axis}
\end{tikzpicture}

\omcaption{Prix illustratifs, non des cotations~: le coût cumulé d'une route longue distance achetée en \omterm{def:nw:long-haul-fibre:dark}{fibre noire} (un IRU de 6~millions et 1.5~million d'équipement
versés d'avance, 0.3~million par an d'exploitation et de maintenance) et louée sous forme de longueur d'onde (90\,000 par mois). Les courbes se croisent après 9.6
ans. Données~: \texttt{fig\_fibre.py}, \texttt{nw\_fibre.costs()}.}\label{fig:nw:long-haul-fibre:costs}
\end{omfigure}

Le calcul du seuil de rentabilité est simple, mais ses paramètres ne le sont pas~: les prix des IRU et les frais des \omterm{def:nw:long-haul-fibre:dark}{services allumés} se négocient et ne sont pas publiés, ce qui explique que la figure utilise
des chiffres illustratifs. Avec eux, la \omterm{def:nw:long-haul-fibre:dark}{fibre noire} coûte 0.675~million par an, amortis sur vingt ans, contre 1.08~million pour la location, et rembourse
son investissement initial au bout de 9.6 ans. Une société qui compte garder la route moins longtemps que cela, ou dont la stratégie pourrait basculer vers le hertzien,
loue~; une société qui a besoin de contrôler chaque microseconde du trajet achète, quelle que soit la réponse du seuil de rentabilité.

\omcode{code/firm/fibreroute/firm_fibreroute.py}{82}{96}{Le modèle de coût de la fibre noire et du service allumé~: coûts annuels et durée du seuil de rentabilité de l'achat
face à la location.}\label{lst:nw:long-haul-fibre:costs}

\section{Verre plus rapide~: fibre à cœur creux}

\begin{definition}[Fibre à cœur creux]\label{def:nw:long-haul-fibre:hcf}
Une \emph{fibre à cœur creux}\index{fibre à cœur creux} (hollow-core fibre) guide la lumière dans un cœur d'air (ou de vide), maintenu par une microstructure de fines membranes de verre
qui l'entoure, de sorte que la lumière se propage presque à $c_0$ au lieu de $c_0/1.46$.
\end{definition}

Pendant des décennies, la \omterm{def:nw:long-haul-fibre:hcf}{fibre à cœur creux} était plus rapide et inutile~: elle perdait trop de lumière pour couvrir de longues distances. Un article de 2025 de l'équipe fibre de Microsoft
et de l'université de Southampton rapporte une \omterm{def:nw:long-haul-fibre:hcf}{fibre à cœur creux} présentant une atténuation de \qty{0.091}{\decibel} par kilomètre à \qty{1550}{\nano\metre},
inférieure à \qty{0.14}{\decibel} de la meilleure fibre en silice massive, avec des débits de transmission supérieurs de 50~\%, «~une latence réduite de 30~\%~» et six fois
moins de \omterm{def:nw:long-haul-fibre:disp}{dispersion chromatique} que la fibre télécom standard. Si l'atténuation se confirme sur les routes câblées, le plancher fibre se rapproche presque du plancher de l'air.

La \Cref{fig:nw:long-haul-fibre:components} reconstruit la route de 2010 de quatre façons avec le modèle (un \omterm{def:nw:the-north-american-map:floor}{indice de groupe} de 1.003 pour le cœur creux est une
hypothèse compatible avec les 30~\% de l'article). Telle qu'inversée, elle prend \qty{6.650}{\milli\second}. Avec une bobine de \omterm{def:nw:long-haul-fibre:disp}{compensation de dispersion} d'un cinquième de
chaque tronçon, elle prendrait \qty{7.976}{\milli\second}~: la compensation optique coûte plus d'une milliseconde sur cette route, d'où toutes les raisons pour une route
à faible latence de compenser électroniquement. Dans la \omterm{def:nw:long-haul-fibre:hcf}{fibre à cœur creux}, sur le même trajet, elle prendrait \qty{4.569}{\milli\second}, \qty{2.08}{\milli\second}
de moins~; dans la fibre standard le long de la géodésique, \qty{5.528}{\milli\second}. Le nouveau verre ferait gagner davantage que n'importe quel redressement du tracé.

\begin{omfigure}
\begin{tikzpicture}
\begin{axis}[xbar stacked,width=11.5cm,height=5.4cm,bar width=10pt,xmin=0,xmax=8.6,ytick=data,
  yticklabels={telle qu'inversée,avec une bobine DCF,cœur creux sur le même tracé,fibre standard sur la géodésique},y dir=reverse,
  xlabel={latence à sens unique, 350 Cermak à Carteret (ms)},tick label style={font=\footnotesize},label style={font=\small},
  grid=major,grid style={gray!25},legend columns=4,legend style={font=\footnotesize,at={(0.5,-0.3)},anchor=north,draw=none,
  /tikz/every even column/.append style={column sep=6pt}},table/col sep=comma]
\addplot[fill=omDef!60,draw=omDef] table[y expr=\coordindex,x=glass]{figdata/networks/13-long-haul-fibre/components.csv};
\addplot[fill=omDef!25,draw=omDef] table[y expr=\coordindex,x=path]{figdata/networks/13-long-haul-fibre/components.csv};
\addplot[fill=omMeth!40,draw=omMeth] table[y expr=\coordindex,x=dcf]{figdata/networks/13-long-haul-fibre/components.csv};
\addplot[fill=omThm!60,draw=omThm] table[y expr=\coordindex,x=equipment]{figdata/networks/13-long-haul-fibre/components.csv};
\legend{verre le long de la géodésique,trajet supplémentaire,DCF,équipement}
\end{axis}
\end{tikzpicture}

\omcaption{La route Chicago--New Jersey de 2010 reconstruite de quatre façons par \texttt{firm.fibreroute}~: par inversion de sa latence cotée, avec compensation optique
de dispersion, en \omterm{def:nw:long-haul-fibre:hcf}{fibre à cœur creux} sur le même trajet, et en fibre standard le long de la géodésique. Les temps d'équipement, la longueur de DCF et
l'indice de la \omterm{def:nw:long-haul-fibre:hcf}{fibre à cœur creux} sont des hypothèses du modèle. Données~: \texttt{fig\_fibre.py}, \texttt{nw\_fibre.variants()}.}\label{fig:nw:long-haul-fibre:components}
\end{omfigure}

\begin{method}[Choisir une route de fibre longue distance]\label{met:nw:long-haul-fibre:choose}
\begin{enumerate}
\item Calculer le plancher entre les deux bâtiments (chapitre~10) et demander à chaque vendeur une latence à sens unique, de baie à baie,
  en précisant comment elle a été mesurée.
\item Inverser chaque chiffre en un \omterm{def:nw:the-north-american-map:factor}{facteur de tracé}, avec une hypothèse d'équipement explicite~; un facteur très supérieur aux autres signifie un long trajet ou
  un équipement lourd, et le vendeur doit indiquer lequel.
\item Demander comment la dispersion est compensée et quelle FEC est utilisée~; demander la diversité de tracé (chapitre~15) si la route doit transporter le trafic
  de plus d'une stratégie.
\item Comparer le prix du dark fibre à celui d'un \omterm{def:nw:long-haul-fibre:dark}{service allumé} sur la durée de vie prévue de la stratégie, et non sur celle de la fibre.
\item Mesurer la route avant de s'y fier, puis une fois en service (chapitres~4 et~5).
\end{enumerate}
\end{method}

\section{Tutoriel~: parcourir une route}\label{tut:nw:long-haul-fibre:tut}

\begin{tutorial}
\textbf{Objectif.} Modéliser une route longue distance comme une suite de tronçons, inverser les latences publiées, et voir ce que changent la \omterm{def:nw:long-haul-fibre:disp}{compensation de dispersion} et la \omterm{def:nw:long-haul-fibre:hcf}{fibre à cœur creux}.
\textbf{État final~:} \cref{fig:nw:long-haul-fibre:sensitivity,fig:nw:long-haul-fibre:costs,fig:nw:long-haul-fibre:components} et
\cref{tab:nw:long-haul-fibre:published}.
\begin{tutsteps}
\item \textbf{Les chiffres publiés.} \texttt{nw\_fibre.PUBLISHED} contient, pour chaque route, ses deux extrémités et sa latence à sens unique, issues du registre~;
  \texttt{published()} les inverse (\cref{lst:nw:long-haul-fibre:components}) avec l'équipement de \texttt{EQUIP}.
\item \textbf{Sensibilité.} \texttt{sensitivity()} répète l'inversion pour des temps d'équipement de 0 à \qty{200}{\micro\second}.
\item \textbf{Variantes.} \texttt{variants()} reconstruit la route avec \texttt{firm\_fibreroute.build} et \texttt{with\_index}.
\item \textbf{Coûts.} \texttt{costs()} et \texttt{break\_even\_years} (\cref{lst:nw:long-haul-fibre:costs}) comparent l'achat et la location
  avec les \texttt{PRICES} illustratifs.
\end{tutsteps}
\textbf{Pour aller plus loin.} Remplacer les hypothèses d'équipement par les chiffres documentés d'un vendeur~; ajouter une portion terrestre à chaque extrémité de la
route atlantique et voir quelle part de son excès ces portions expliquent.
\end{tutorial}

\section{Construction~: le modèle de route de fibre}\label{bld:nw:long-haul-fibre:fibreroute}

\begin{build}
\bfield{Rôle} Routes de fibre longue distance sous forme de données et de calculs~: leur latence par composant, ce qu'implique un chiffre publié, et ce que coûtent leur achat ou
leur location~; utilisé par les chapitres~14 (la fibre de secours d'une liaison radio), 15 (l'achat) et 29 (le plan).
\bfield{Interface} \texttt{firm\_fibreroute}~: \texttt{Span(km, n\_g, amp\_us, dcf\_km)}, \texttt{Route(name, geodesic\_km, spans, terminal\_us)},
\texttt{build}, \texttt{components}, \texttt{invert}, \texttt{with\_index}, \texttt{dark}, \texttt{lit}, \texttt{break\_even\_years}.
\bfield{Règles} Un trajet n'est jamais plus court que sa géodésique~; chaque temps d'équipement est un paramètre explicite~; les latences publiées sont à sens unique et
  indiquent comment elles ont été obtenues.
\bfield{Tests d'acceptation} \texttt{code/firm/fibreroute/tests/}~: les composants à la main, l'aller-retour de l'inversion, le rapport du cœur creux, une bobine de DCF, et
l'arithmétique du point d'équilibre, y compris un leasing toujours moins cher.
\bfield{Pour aller plus loin} Routes sous forme de polylignes sur les coordonnées de \texttt{firm.geomap}~; perte et placement des amplificateurs à partir de l'atténuation.
\end{build}

\begin{omsources}
\item G.~Laughlin, A.~Aguirre et J.~Grundfest, ``Information transmission between financial markets in Chicago and New York'', \emph{Financial
Review} 49(2) (2014).
\item M.~Petrovich, E.~Numkam Fokoua et al., ``Broadband optical fibre with an attenuation lower than 0.1 decibel per kilometre'', \emph{Nature
Photonics} (2025).
\item Corning, SMF-28 ULL product information (2020)~; Equinix, communiqué de presse sur Hibernia Express (3 novembre 2015)~; Cogent Communications, Form 10-K
pour 2022.
\end{omsources}

\section{Exercices}

\begin{exercise}[$\star$]\label{exo:nw:long-haul-fibre:1}
Avec une atténuation de \qty{0.16}{\decibel} par kilomètre, quelle fraction de la lumière subsiste après un tronçon de 80~kilomètres~?
\end{exercise}

\begin{exercise}[$\star$]\label{exo:nw:long-haul-fibre:2}
En utilisant la \cref{prop:nw:long-haul-fibre:invert}, quelle quantité de verre contient une route publiée à \qty{6.55}{\milli\second} si son équipement prend
\qty{21.7}{\micro\second}~?
\end{exercise}

\begin{exercise}[$\star$]\label{exo:nw:long-haul-fibre:3}
Quelle est la différence entre le dark fibre et un \omterm{def:nw:long-haul-fibre:dark}{service de longueur d'onde}, du point de vue de la latence~?
\end{exercise}

\begin{exercise}[$\star\star$]\label{exo:nw:long-haul-fibre:4}
Pourquoi la compensation optique de dispersion coûte-t-elle de la latence, et la compensation électronique non (ou beaucoup moins)~?
\end{exercise}

\begin{exercise}[$\star\star$]\label{exo:nw:long-haul-fibre:5}
Le \omterm{def:nw:the-north-american-map:factor}{facteur de tracé} de la route atlantique vaut 1.09 et celui de la route de Chicago 1.20. Donner deux raisons.
\end{exercise}

\begin{exercise}[$\star\star$]\label{exo:nw:long-haul-fibre:6}
Un vendeur annonce un aller-retour. Que faut-il demander avant de le diviser par deux~?
\end{exercise}

\begin{exercise}[$\star\star\star$]\label{exo:nw:long-haul-fibre:7}
\emph{Programmation.} Avec \texttt{firm.fibreroute}, trouver le \omterm{def:nw:the-north-american-map:factor}{facteur de tracé} auquel une route en \omterm{def:nw:long-haul-fibre:hcf}{fibre à cœur creux} entre Cermak et Carteret égalerait une
route hertzienne de \omterm{def:nw:the-north-american-map:factor}{facteur de tracé} 1.011 dans l'air.
\end{exercise}

\begin{exercise}[$\star\star\star$]\label{exo:nw:long-haul-fibre:8}
\emph{Trouver la faille.} «~Notre nouvelle route en fibre n'est que 3~\% plus longue que la géodésique, elle battra donc les réseaux hertziens, dont les pylônes zigzaguent.~»
\end{exercise}

\section{Problème~: Treize millisecondes}

\begin{problem}[{Problème du week-end --- démonter la route Chicago--New Jersey de 2010}]\label{pb:nw:long-haul-fibre:1}
La route de fibre de 2010 entre 350 Cermak et Carteret était cotée à \qty{6.65}{\milli\second} à sens unique. Utiliser les hypothèses du modèle~: \qty{10}{\micro\second}
par terminal, un amplificateur tous les 80~kilomètres à \qty{0.1}{\micro\second}, et une fibre standard d'\omterm{def:nw:the-north-american-map:floor}{indice de groupe} 1.462.

\medskip\noindent\textbf{Partie I --- Le plancher.}
\begin{enumerate}
\item Quelles sont la \omterm{def:nw:the-north-american-map:geodesic}{distance géodésique} et les planchers dans le vide et dans la fibre entre les deux bâtiments~?
\item Quelle est la latence aller-retour cotée, et le plancher aller-retour dans la fibre~?
\item Quelle part du chiffre à sens unique dépasse le plancher de la fibre~?
\item Pourquoi le plancher de la fibre, et non celui du vide, est-il la bonne comparaison pour une route en fibre~?
\end{enumerate}

\medskip\noindent\textbf{Partie II --- La répartition.}
\begin{enumerate}[resume]
\item Quelle quantité de verre contient la route, et combien de tronçons~?
\item Quel est le temps d'équipement dans le modèle, et quelle part de l'excès est due à l'équipement~?
\item Quelle part représente le trajet supplémentaire, en microsecondes et en kilomètres~?
\item Que devient la répartition si l'équipement prend \qty{200}{\micro\second}~?
\end{enumerate}

\medskip\noindent\textbf{Partie III --- Autre verre.}
\begin{enumerate}[resume]
\item Que donnerait une \omterm{def:nw:long-haul-fibre:disp}{compensation de dispersion} optique sur un cinquième de chaque tronçon~?
\item Combien prendrait le même trajet en \omterm{def:nw:long-haul-fibre:hcf}{fibre à cœur creux} d'\omterm{def:nw:the-north-american-map:floor}{indice de groupe} 1.003~?
\item Et en fibre standard le long de la géodésique~?
\item Comment ces valeurs se comparent-elles à la route hertzienne de 2016 du chapitre~10 (\qty{3.982}{\milli\second} depuis Aurora)~?
\end{enumerate}

\medskip\noindent\textbf{Partie IV --- Le verdict.}
\begin{enumerate}[resume]
\item Énoncer le \emph{résultat nommé}~: la décomposition de l'aller-retour de 2010 en verre le long de la géodésique, trajet supplémentaire et équipement, ainsi que les millisecondes
  qu'une \omterm{def:nw:long-haul-fibre:hcf}{fibre à cœur creux} sur le même trajet permettrait de gagner.
\item Quelle partie de la décomposition dépend le plus des hypothèses du modèle~?
\item Pourquoi la route a-t-elle perdu son avance en moins de deux ans~?
\item Qu'a-t-elle conservé~?
\item L'achèteriez-vous en \omterm{def:nw:long-haul-fibre:dark}{fibre noire} ou loueriez-vous une longueur d'onde sur celle-ci~?
\item Que faudrait-il à une route à cœur creux pour battre les réseaux hertziens~?
\item Que dit le \omterm{def:nw:the-north-american-map:factor}{facteur de tracé} de la route atlantique sur l'endroit où la route de Chicago a perdu son temps~?
\item En une phrase~: où passent les treize millisecondes~?
\end{enumerate}
\end{problem}

\section{Questions d'entretien}

\begin{interviewq}[$\star$ \iqroles{développeur}]\label{iq:nw:long-haul-fibre:1}
À quelle vitesse la lumière se propage-t-elle dans la fibre, et quelle est la latence par tranche de 1\,000 kilomètres~?
\end{interviewq}

\begin{interviewq}[$\star\star$ \iqroles{développeur}]\label{iq:nw:long-haul-fibre:2}
Qu'est-ce qui ajoute de la latence à une route de fibre longue distance, en dehors de la longueur du verre~?
\end{interviewq}

\begin{interviewq}[$\star\star$ \iqroles{développeur, trader}]\label{iq:nw:long-haul-fibre:3}
Pourquoi les faisceaux hertziens ont-ils battu la fibre entre Chicago et New Jersey, et pourquoi la fibre est-elle encore utilisée là-bas~?
\end{interviewq}

\begin{interviewq}[$\star\star$ \iqroles{développeur}]\label{iq:nw:long-haul-fibre:4}
Qu'est-ce qu'un IRU, et quand une société de trading achèterait-elle de la \omterm{def:nw:long-haul-fibre:dark}{fibre noire} plutôt que de louer une longueur d'onde~?
\end{interviewq}

\begin{interviewq}[$\star\star$ \iqroles{développeur, chercheur}]\label{iq:nw:long-haul-fibre:5}
Qu'est-ce que la \omterm{def:nw:long-haul-fibre:hcf}{fibre à cœur creux}, et que changerait-elle pour le trading sensible à la latence~?
\end{interviewq}

\begin{interviewq}[$\star\star\star$ \iqroles{développeur}]\label{iq:nw:long-haul-fibre:6}
Un vendeur affirme qu'une route de fibre entre deux villes est plus rapide que votre plancher de fibre calculé. Qu'est-ce qui pourrait l'expliquer~?
\end{interviewq}
